\chapter{Pretraining}

\section{Cross-Entropy Loss for Next-Token Prediction}

The model outputs logits of shape $[\text{batch} \times \text{seq\_len},
\text{vocab\_size}]$. The target is the shifted token IDs of shape
$[\text{batch} \times \text{seq\_len}]$.

\[
  \mathcal{L} = -\frac{1}{N}\sum_{i=1}^{N} \log P(w_{i+1} \mid w_1, \ldots, w_i)
\]

\begin{keyconcept}
  Cross-entropy expects:
  \begin{itemize}
    \item \textbf{Input}: logits $[\text{batch} \times \text{seq}, \text{vocab}]$
    \item \textbf{Target}: integer class indices $[\text{batch} \times \text{seq}]$
  \end{itemize}
  The target is the input shifted by 1 position: ``The cat sat'' $\rightarrow$
  ``cat sat on''.
\end{keyconcept}

\section{Training Loop with AdamW}

\begin{pseudocode}[Training Loop]
\textbf{for} epoch $= 1$ \textbf{to} $E$:\\
\quad \textbf{for} batch \textbf{in} dataloader:\\
\quad\quad logits $\leftarrow$ model(batch.input)\\
\quad\quad loss $\leftarrow$ CrossEntropy(logits, batch.target)\\
\quad\quad loss.backward()\\
\quad\quad optimizer.step()\\
\quad\quad optimizer.zero\_grad()
\end{pseudocode}

\begin{keyconcept}
  \textbf{AdamW} decouples weight decay from the gradient update:
  \[
    \theta_{t+1} = \theta_t - \eta\left(\frac{\hat{m}_t}{\sqrt{\hat{v}_t} + \epsilon} + \lambda \theta_t\right)
  \]
  Typical hyperparameters: $\text{lr} = 3 \times 10^{-4}$,
  weight decay $= 0.01$, $\beta_1 = 0.9$, $\beta_2 = 0.999$.
\end{keyconcept}

\section{Learning Rate Scheduling}

The learning rate follows a \textbf{warmup + cosine decay} schedule:

\[
  \text{lr}(t) =
  \begin{cases}
    \text{max\_lr} \cdot \frac{t}{T_w} & \text{if } t \leq T_w \\[6pt]
    \text{min\_lr} + \frac{1}{2}(\text{max\_lr} - \text{min\_lr})
    \left(1 + \cos\!\left(\pi \cdot \frac{t - T_w}{T - T_w}\right)\right)
    & \text{if } t > T_w
  \end{cases}
\]

where $T_w$ is the number of warmup steps (typically 10\% of total) and
$T$ is the total number of training steps. $\text{min\_lr}$ is typically
$\text{max\_lr} / 10$.

\begin{keyconcept}
  \textbf{Warmup} prevents instability at the start of training when
  gradients are large and the model is randomly initialized. \textbf{Cosine
  decay} allows the learning rate to decrease smoothly, enabling
  finer parameter adjustments in later stages of training.
\end{keyconcept}

\section{Weight Initialization}

GPT-2 initializes weights from $\mathcal{N}(0, 0.02)$ for embeddings
and linear layers. This small standard deviation prevents early
training instability.

\section{Gradient Accumulation}

When batch size is limited by GPU memory, gradient accumulation
simulates larger batches:

\begin{pseudocode}[Gradient Accumulation]
\textbf{for} step $= 1$ \textbf{to} total\_steps:\\
\quad batch\_loss $\leftarrow$ model(batch)\\
\quad scaled\_loss $\leftarrow$ batch\_loss / grad\_accml\_steps\\
\quad scaled\_loss.backward() \hfill $\triangleright$ Accumulate gradients\\
\quad \textbf{if} step \% grad\_accml\_steps $== 0$:\\
\quad\quad optimizer.step()\\
\quad\quad optimizer.zero\_grad()
\end{pseudocode}

\begin{keyconcept}
  Dividing the loss by \texttt{grad\_accml\_steps} ensures that the
  effective gradient magnitude is the same as with a true large batch.
  Without this division, gradients would be too large by a factor of
  \texttt{grad\_accml\_steps}.
\end{keyconcept}

\section{Mixed Precision Training}

Using \texttt{torch.autocast} with \textbf{bfloat16} (or float16):

\begin{pseudocode}[Mixed Precision]
\textbf{with} autocast(device\_type="cuda", dtype=bfloat16):\\
\quad logits $\leftarrow$ model(batch)\\
\quad loss $\leftarrow$ criterion(logits, targets)\\[3pt]
loss.backward()\\
optimizer.step()
\end{pseudocode}

\begin{keyconcept}
  Mixed precision stores weights in float32 but computes forward/backward
  passes in bfloat16, halving memory usage and often doubling throughput
  on modern GPUs. Bfloat16 has the same dynamic range as float32 (8 exponent
  bits), avoiding the overflow/underflow issues of float16.
\end{keyconcept}

\section*{Review Questions}

\begin{reviewquestion}[1]
  Why does gradient accumulation divide the loss by the number of
  accumulation steps? What would happen to the effective learning rate
  without this division?
\end{reviewquestion}

\begin{reviewquestion}[2]
  Sketch the warmup + cosine decay learning rate curve. What happens if
  you skip warmup and start at max\_lr immediately?
\end{reviewquestion}

\begin{reviewquestion}[3]
  Why does bfloat16 have a wider dynamic range than float16? How many
  exponent bits does each format use?
\end{reviewquestion}